\documentclass[11pt]{article}
\usepackage[T1]{fontenc}
\usepackage{lmodern,amsmath,amssymb,amsthm,mathtools,booktabs,enumitem,microtype}
\usepackage[margin=1in]{geometry}
\usepackage[hidelinks]{hyperref}
\usepackage[capitalise,noabbrev]{cleveref}
\allowdisplaybreaks[2]
\setlength{\emergencystretch}{1.6em}
\setlist{itemsep=.25em,topsep=.45em}
\newtheorem{theorem}{Theorem}
\newtheorem{lemma}[theorem]{Lemma}
\newtheorem{proposition}[theorem]{Proposition}
\newtheorem{corollary}[theorem]{Corollary}
\newtheorem{claim}[theorem]{Claim}
\newtheorem{fact}[theorem]{Fact}
\theoremstyle{definition}
\newtheorem{definition}[theorem]{Definition}
\newtheorem*{remark}{Remark}
\newcommand{\R}{\mathbb R}
\newcommand{\N}{\mathbb N}
\newcommand{\abs}[1]{\lvert#1\rvert}
\newcommand{\norm}[1]{\lVert#1\rVert}
\newcommand{\pos}[1]{(#1)_+}
\newcommand{\inner}[2]{\langle #1,#2\rangle}
\newcommand{\psil}{p_{\mathrm{sil}}}
\DeclareMathOperator{\prox}{prox}
\DeclareMathOperator{\conv}{conv}
\DeclareMathOperator{\dist}{dist}
\DeclareMathOperator{\supp}{supp}
\DeclareMathOperator{\env}{env}
\hypersetup{pdftitle={\detokenize{Sharp Omega(kappa^(1/p_sil) log(1/epsilon)) Lower Bound for Gradient Descent on Strongly Convex Objectives}},pdfauthor={Yuhan Ye}}
\title{Sharp $\Omega\!\left(\kappa^{1/\psil}\log(1/\varepsilon)\right)$ Lower Bound for Gradient Descent\\on Strongly Convex Objectives}
\author{Yuhan Ye}
\date{November 25, 2026}
\makeatletter
\renewcommand{\maketitle}{\begin{center}{\LARGE\@title\par}\medskip\@author\par\smallskip\@date\end{center}}
\makeatother
\begin{document}
\maketitle
\begin{abstract}
Altschuler and Parrilo~\cite{AP} showed that their \emph{silver stepsize schedule} accelerates gradient descent (GD) on smooth strongly convex functions. To halve the error, it improves the classical $O(\kappa)$ iteration bound to $O(\kappa^{1/\psil})$, where $\kappa$ is the condition number and $\psil=\log_2(1+\sqrt2)$. For nonnegative stepsizes, Liu and Ye~\cite{LY} matched this bound up to a subpolynomial factor in $\kappa$.

In this paper, we prove that GD with any predetermined sequence of real stepsizes requires $\Omega(\kappa^{1/\psil}\log(1/\varepsilon))$ iterations in the worst case to reduce either the squared distance to the minimizer or function suboptimality by a factor $\varepsilon$. For $0<\varepsilon\le1/2$, this matches the silver upper bound up to an absolute constant.
\end{abstract}

\section{Introduction}
How much can the choice of stepsizes accelerate gradient descent? On an $L$-smooth, $\mu$-strongly convex function $F$ with minimizer $x_*$, GD with stepsizes $h_1,\ldots,h_n$ takes the steps
\[
 x_t=x_{t-1}-h_t\nabla F(x_{t-1}).
\]
We consider \emph{predetermined} stepsizes, chosen before the run. The classical choice $h_t=1/L$ takes $O(\kappa)$ iterations to halve the squared distance to $x_*$, where $\kappa=L/\mu$. Altschuler and Parrilo~\cite{AP} reduced this bound to $O(\kappa^q)$ by combining short and long steps in their silver schedule, where $q=1/\psil$ and $\psil=\log_2(1+\sqrt2)$.

For nonnegative stepsizes, Liu and Ye~\cite{LY} proved that halving the error requires $\Omega(\kappa^{q-o(1)})$ iterations. This leaves open whether the subpolynomial gap can be closed and whether negative stepsizes can improve the rate. Jeong and Hong~\cite{JH} recently proved that the silver rate is sharp for real stepsizes on smooth convex objectives. We prove the corresponding result for strongly convex objectives.

\paragraph{Main result.}
We normalize $L=1$, so that $\mu=1/\kappa$. Throughout, $\log$ denotes the natural logarithm, and $c,C$ denote positive absolute constants unless a dependence is specified. For a predetermined schedule $H=(h_1,\ldots,h_n)\in\R^n$, define the worst-case ratios for squared distance and function suboptimality by
\[
 D(H)=\sup_{F,x_0}\frac{\|x_n-x_*\|^2}{\|x_0-x_*\|^2},
 \qquad E(H)=\sup_{F,x_0}\frac{F(x_n)-F(x_*)}{F(x_0)-F(x_*)}.
\]
Each supremum ranges over all finite dimensions, all $1$-smooth, $\mu$-strongly convex objectives $F$, and all initial points $x_0\ne x_*$.

\begin{theorem}\label{thm:main}
There exist absolute constants $C,\kappa_0>0$ such that for every $\kappa\ge\kappa_0$, every $n\ge1$, and every $H\in\R^n$,
\[
 \min\{D(H),E(H)\}\ge \exp(-Cn/\kappa^q).
\]
Thus a uniform reduction of either error by a factor $\varepsilon\in(0,1)$ requires $\Omega(\kappa^q\log(1/\varepsilon))$ steps.
\end{theorem}

To see the matching upper bound for both errors, take a silver block of length $s=2^k-1$. Its guarantee $F(x_s)-F(x_*)\le CL\|x_0-x_*\|^2/s^{\psil}$~\cite[Theorem~1]{APconv} and strong convexity show that a block of length $s\asymp\kappa^q$ with a sufficiently large constant halves both errors. Repeating it gives the stated order for $0<\varepsilon\le1/2$.

\paragraph{Proof idea.}
In the construction of Liu and Ye~\cite[Sections~4--5]{LY}, selected positive steps, called \emph{checkpoints}, activate successive coordinates of the hard function.\footnote{For smooth convex objectives, Ma and Chen~\cite[Theorem~4.1]{MC} chose long steps at which the gradient changes direction. Jung, Cho, and Yun~\cite[Section~2]{JCY} realized these changes with local functions coupling adjacent coordinates.} Between checkpoints, the quadratic term that ensures strong convexity can shrink the active coordinate. We therefore keep track of the total length of the short positive steps and the number of long ones. Fix $K=\kappa^{1/1000}$ and write
\[
 S_K^+(H)=\sum_{t=1}^n\min\{(h_t)_+,K\},\qquad (h)_+=\max\{h,0\}.
\]
Lemma~\ref{lem:block} constructs a hard function when this sum is small enough. The difficulty with signed steps is that a negative step can bring the iterate back to an earlier point. Where the chain cannot pass to a new coordinate, we use three quadratic coordinates of different curvatures. Together they distinguish the iterate at the next checkpoint from earlier iterates, except for visits separated by steps of very small total absolute length. This lets us start the next part of the chain on the same function.

Lemma~\ref{lem:count} proves the required bound on $S_K^+(H)$. We adapt Jeong and Hong's argument for smooth convex objectives~\cite[Sections~3--7]{JH}. They count steps above a varying threshold and use explicit GD trajectories to bound this count. In the strongly convex setting we apply the count on shorter intervals and join their trajectories on one objective. A tool for this last step comes from the performance estimation framework of Drori and Teboulle~\cite{DT}. The interpolation criterion of Taylor, Hendrickx, and Glineur~\cite[Corollary~1]{THG} characterizes which finite lists of values and gradients come from one smooth convex function. At a switch we use these inequalities to replace part of the objective by a quadratic on a small ball while preserving its values and gradients outside a larger ball (Appendix~\ref{app:arc}).

\begin{lemma}[Constructing the hard function]\label{lem:block}
Fix $C_*>0$. If $n\ge\kappa^q$ and $S_K^+(H)\le C_*nK^{1-q}$, then for all sufficiently large $\kappa$ there are a $1$-smooth, $\mu$-strongly convex function $F$ and an initial point $x_0\ne x_*$ with
\[
 \|x_n-x_*\|/\|x_0-x_*\|\ge \exp(-A n/\kappa^q),
\]
where $A$ and the lower cutoff on $\kappa$ depend only on $C_*$.
\end{lemma}

\begin{lemma}[Bounding the positive steps]\label{lem:count}
For some absolute $C_*$ and all sufficiently large $\kappa$, if $n\ge\kappa^q$ and $\min\{D(H),E(H)\}\le1$, then $S_K^+(H)\le C_*nK^{1-q}$.
\end{lemma}

\begin{proof}[Proof of Theorem~\ref{thm:main}]
It suffices to consider $\min\{D(H),E(H)\}\le1$. Repeat $H$ a total of $r$ times. The resulting schedule $H^r$ satisfies $D(H^r)\le D(H)^r$ and $E(H^r)\le E(H)^r$. For $rn\ge\kappa^q$, Lemmas~\ref{lem:count} and~\ref{lem:block} give $D(H^r)\ge e^{-2Arn/\kappa^q}$. Strong convexity at $x_{rn}$ and smoothness at $x_0$ give $E(H^r)\ge D(H^r)/\kappa$. Hence
\[
 \min\{D(H),E(H)\}\ge \kappa^{-1/r}\exp(-2An/\kappa^q).
\]
Letting $r\to\infty$ proves the theorem for every $n$.
\end{proof}

\section{Proof of Lemma~\ref{lem:block}}\label{sec:proof-block}

We adapt the blockwise checkpoint selection of Liu and Ye~\cite[Sections~4--5]{LY}.
We select checkpoints using the positive parts \(p_t=(h_t)_+\) and run GD with
the original stepsizes \(h_t\). If the sum of the unselected positive stepsizes
becomes too large, we add a checkpoint, a procedure called a \emph{repair}.
Usually this checkpoint activates a new coordinate. When the step it selects
is too short, we use three quadratic coordinates until the chain can resume.

\paragraph{Transfer at a checkpoint.}
Call the sum of \(p_t\) over an interval its \emph{positive mass}.
Consider a checkpoint of stepsize \(b\). Its \emph{gap} consists of the steps
after the preceding checkpoint (or from initialization) and strictly before
this one. Let \(S\) be the positive mass of the gap. At its start, the input
coordinate exceeds its activation threshold by \(D>0\). Suppose that threshold
is at most \(D/2\), \(S\le\delta\kappa\) for a sufficiently small absolute
constant \(\delta>0\), and \(b>K\). The construction in Appendix~\ref{app:arc}
makes the output coordinate exceed its threshold by \(D_+>0\) at the checkpoint, with
\begin{equation}\label{eq:signed-block-transfer}
 \frac{D_+}{D}\ge\frac{b-8}{c_K(K+S+8)}e^{-C\mu(K+S)},
 \qquad 1\le c_K\le1+C/\log K.
\end{equation}
The coefficient \(c_K\) depends only on \(K\). The outgoing threshold is at most \(D_+/2\), as required for the next gap.
The dependence on positive mass comes from the quadratic term. For
\(\lambda\ge0\) with \(\lambda S<1\), every prefix of the gap satisfies
\(\prod_t(1-\lambda h_t)\ge1-\lambda S\), since each negative step has
multiplier at least one. The component is \(4\mu\)-strongly convex and
\(1/4\)-smooth, leaving room to join trajectories later in the proof.

\paragraph{Selecting checkpoints.}
Put \(\nu=1/\log_2(1+\sqrt{1+c_K})\). Since
\(0\le q-\nu=O(1/\log K)\), we have \(\kappa^\nu\asymp\kappa^q\).
Choose a small absolute \(\tau>0\) below, and let
\(U=(\tau/2)^{1/\nu}\kappa\). Split each run of steps \(p_t\le K\)
greedily into pieces of positive mass at most \(U/2\). Represent a piece
of mass \(s\) immediately followed by a large step \(p_t>K\) by the pair
\((g,b)=(s+8,p_t-8)\). An initial large step or a large step immediately
following another has \(s=0\). A piece without a following large step
gives the pair \((s,0)\). In each case \(g+b\) equals the represented mass.

The bound on \(S_K^+(H)\) in Lemma~\ref{lem:block} allows \(O(nK^{-q})\) large steps
and \(O(nK^{1-q})\) mass in the other steps. Each additional cut
accounts for mass at least \(U/2-K\). Since \(n\ge\kappa^q\), the
number \(N\) of pairs is \(O(nK^{-q})\) and \(\sum g=O(nK^{1-q})\).
By concavity,
\(\sum(K+g)^\nu\le N^{1-q}(KN+\sum g)^q=O(n)\).
Group consecutive pairs greedily into blocks of weight
\(\sum(K+g)^\nu\le\tau\kappa^\nu\). Each pair uses at most half
this allowance, so the number \(B\) of blocks is \(O(n/\kappa^q)\).

At the start of a block, let \(w\) be the positive mass after the latest
checkpoint and before this block, including all earlier steps if no
checkpoint has been selected. Add \(w\) to the first \(g\). Among all sets
of indices \(t_1<\cdots<t_k\) in this block with \(h_{t_i}>K\), maximize
\(\prod_i b_i/d_i\), where \(b_i=h_{t_i}-8\), \(S_i\) is the positive
mass of the gap ending at \(t_i\), and \(d_i=c_K(K+S_i+8)\). For the first
gap \(S_1\) includes \(w\). The empty product is one. In the alternating
list of paired entries \((g,b)\), this is the score of~\cite[Appendix~B]{LY}
with offset \(K\). Let \(M=w^\nu+\tau\kappa^\nu\), an upper
bound on the weight including \(w\). We use three consequences of
\cite[Lemmas~B.5--B.7 and proof of Lemma~4.1(ii)]{LY}. When the
maximizer is empty the block including its incoming mass has total
mass \(O(M^{1/\nu})\) and at most \(c_K^\nu M x^{-\nu}\) entries
\(b>x\) for every \(x>0\). Otherwise the selected denominators \(d_i\)
sum to at most \(c_K M^{1/\nu}\), and the suffix after the last
selection has an empty maximizer with zero incoming mass.

Choose \(C_b\) so these mass bounds are at most
\(C_b(w+\tau^{1/\nu}\kappa)\). Set \(S_0=\delta\kappa/(4C_b)\)
and choose \(\tau\) so \(C_b\tau^{1/\nu}\kappa\le S_0/4\).
A nonempty maximizer leaves at most \(S_0/4\) unselected mass.
If the maximizing set is empty, the \emph{unselected tail} consists of all
steps since the latest checkpoint through the end of the current block.
Carry this tail forward unless its mass reaches \(S_0\). In that case take its shortest suffix of mass
greater than \(S_0/8\) and choose a largest step in it as a repair
checkpoint. This leaves at most \(S_0/8\) mass after the repair.
Consequently \(w<S_0\) throughout. Both the sum of the selected
denominators in a block and the mass of a tail presented for repair
are at most \(2C_b S_0\). In particular every gap and final tail has
positive mass at most \(\delta\kappa\).

\paragraph{Bounding the loss at repair checkpoints.}
For repair \(j\), let \(I_j\) be its shortest suffix of positive mass
greater than \(S_0/8\), and let \(h_j^\star=\max_{t\in I_j}p_t\)
be the stepsize selected for repair. These suffixes are disjoint.
Indeed, at most \(S_0/8\) mass remains at the end of a repair block,
so the next repair needs at least \(7S_0/8\) new mass. Its shortest
suffix therefore begins after that block. If \(I_j\) meets \(L_j\)
blocks, it crosses \(L_j-1\) block boundaries. Hence
\(\sum_j L_j\le 2B\), since there is at most one repair per block.

Call a repair exceptional if the steps \(p_t\le K\) in \(I_j\)
have total mass greater than \(S_0/16\). By disjointness and the
bound on \(S_K^+(H)\), their number \(J\) satisfies
\begin{equation}\label{eq:signed-exception-count}
 J\le \frac{16}{S_0}\sum_{t:p_t\le K}p_t
 \le C(n/\kappa^q)\kappa^{-(1-q)(1-1/1000)} .
\end{equation}

In every other repair the steps larger than \(K\) in \(I_j\)
have total mass greater than \(S_0/16\). We claim that
\(h_j^\star/S_0\ge c_1 L_j^{-1/(1-q)}\). This holds immediately
if \(h_j^\star\ge S_0\). Otherwise, write \(N_j(x)\) for the
number of steps \(p_t>x\) in \(I_j\). Each block contributes only
steps from a block whose maximizing set is empty or from the suffix after
the last selected checkpoint. For \(K\le x\le S_0\), every \(p_t>x\)
has encoded size \(p_t-8>x/2\). The level count above, with
\(M=O(S_0^\nu)\) and \(\nu\le q\), gives
\(N_j(x)\le C L_j(S_0/x)^q\). Consequently
\[
 \frac{S_0}{16}
 < K N_j(K)+\int_K^{h_j^\star} N_j(x)\,dx
 \le C L_j S_0^q(h_j^\star)^{1-q},
\]
which proves the claim. These repairs all select \(h_j^\star>K\).
Their entire incoming gap has positive mass \(S\le 2C_bS_0\), so
\(K+S+8=O(S_0)\), \(\mu(K+S)=O(1)\), and \(h_j^\star-8\ge
h_j^\star/2\) for large \(\kappa\). By \eqref{eq:signed-block-transfer},
\(D_+/D\ge c_2 h_j^\star/S_0\). Summed over these repairs, the
logarithmic loss is at most \(O(\sum_j(1+\log L_j))=O(B)\).

\paragraph{Connecting the exceptional intervals.}
A nonexceptional checkpoint is available for changing coordinates if
its size lies in \([4,\kappa/4]\). All nonexceptional repairs have
this property. Include \(0,n+1\) as endpoints, and mark every open
interval between successive available indices that contains an
exceptional repair. There are at most \(J\) marked intervals.
At the entrance and exit of each, reserve the boundary checkpoints
for changing trajectories. A shared boundary is used just once.

In an unmarked portion use the components from \eqref{eq:signed-block-transfer}.
A whole maximizing set has score
at least one, and \(\sum d_i=O(\kappa)\) bounds its exponential loss by
a constant in the logarithm. A cut can leave a contiguous fragment
\(I\) of that set. Its log loss \(A(I)=\sum_{i\in I}\log(d_i/b_i)\)
is \(O(\log\kappa\log(1+|I|))\). Indeed, let \(b_*=\max_{i\in I}b_i\).
Deleting the fragment changes only the next surviving denominator,
increasing it by at most \(C(\kappa+|I|(b_*+8))\), so maximality
gives \(A(I)\le\log b_*+O(\log\kappa)\) (or \(A(I)\le0\) for a
terminal fragment). Here \(|I|\le C\kappa\). Let \(I_1,I_2\) be the
subsets before and after an index attaining \(b_*\). The contribution at
that index is at most \(O(\log\kappa)-\log b_*\), so
\(2A(I)\le A(I_1)+A(I_2)+O(\log\kappa)\). Induction proves the
claim since the average of \(\log(1+|I_1|),\log(1+|I_2|)\) is
at most \(\log(1+|I|)-\log2\). There are \(O(J)\) cut fragments.

On a marked interval every positive step is either at most
\(\delta\kappa\) or greater than \(\kappa/4\): the open gaps and
exceptional repair tails have mass at most \(\delta\kappa\), and any
other checkpoint there exceeds \(\kappa/4\). Use three quadratics
with curvatures \(12\mu,16\mu,20\mu\). On each of them a large step
multiplies the distance by at least two and pays for the log loss
\(O(\delta)\) in the gap preceding it. Negative steps only increase
the distance. Exceptional checkpoints, their preceding gaps, and the
final gap cost \(O(1+J_I)\) in log distance, where \(J_I\) is the
number of exceptional repairs in that interval.

Appendix~\ref{app:arc} proves that these portions can be realized on
one objective. Here is the quantitative estimate needed to finish.
Start each portion at unit norm. For an arc portion let \(G_i\) be
its final excess, just before the last open gap. The output stays
at least \(c G_i\) through this gap since its positive mass is at most
\(\delta\kappa\) and its incoming threshold is at most \(G_i/2\);
at initialization take \(G_i=1\). For a quadratic portion let
\(G_i\) be the smallest of its three final absolute multipliers.
The estimates above give
\(\sum_i\log^-(G_i)=O(B+J\log^2\kappa)\), where
\(\log^-x=\max(0,-\log x)\).

At a reserved step following a portion of length \(r_i\), the three
curvatures distinguish its last query from earlier ones. Each nonzero
step increases the middle multiplier's magnitude relative to
the geometric mean of the two outer magnitudes. The construction giving
\eqref{eq:joins-gain} then starts the next portion exactly at a scaled
copy of its prescribed input. In addition to
\(\log^-(G_i)\), the log cost per switch is at most
\(C(1+\log\kappa+\log(1+r_i)+\log^-(G_{i+1}))\).
There are at most \(2J\) switches, and \(\sum r_i\le n\).
The last portion needs no switch. Thus the total loss is at most
\[
 C\{B+J\log^2\kappa+J\log(1+n/J)\}=O(n/\kappa^q),
\]
where the last term is zero if \(J=0\). The equality follows from
\eqref{eq:signed-exception-count} and monotonicity of
\(x\log(1+n/x)\). Exponentiating proves Lemma~\ref{lem:block}.
\section{Proof of Lemma~\ref{lem:count}}\label{sec:proof-count}

We bound \(S_K^+(H)\) by counting the stepsizes above each threshold.
Jeong and Hong~\cite[Sections~3--7]{JH} obtained such a count in the
smooth convex setting. At each threshold they select the steps above it
as checkpoints and construct a function on which the ratio of final to
initial error depends on those checkpoints. Here we first cut the
schedule into intervals on which the quadratic contraction is small.
We adapt their count on each interval, then join the corresponding
trajectories on one strongly convex function.

\paragraph{Cutting at large steps.}
On \(f(x)=\lambda x^2/2\), both error ratios equal \(|P(\lambda)|^2\)
for \(P(z)=\prod_{j=1}^n(1-h_jz)\). Since the schedule is stable in
at least one of the two measures, that is, \(\min\{D(H),E(H)\}\le1\), we have \(|P|\le1\) on
\([\mu,1]\). For \(1\le r\le\kappa\) it follows that
\begin{equation}\label{eq:count-quadratic}
 \sum_j\log(1+h_j^2/r^2)=2\log|P(\mathrm i/r)|\le Cn/\sqrt r .
\end{equation}
Indeed the affine image of \(\mathrm i/r\) under
\([\mu,1]\mapsto[-1,1]\) is \(z=-1+O(1/r)\). Writing
\(z=(w+w^{-1})/2\) with \(|w|\ge1\) gives \(|w|\le1+C/\sqrt r\).
The maximum principle on the exterior disk, applied to the transformed
polynomial divided by \(w^n\), gives \(|P(\mathrm i/r)|\le |w|^n\).

Set \(\hat h_j=h_j/4\). A step of size \(h_j\) on \(F/4\) then gives
the same iterate as a step of size \(\hat h_j\) on \(F\). This scaling
leaves room to join consecutive intervals while preserving the required
smoothness and strong convexity. Remove all steps with \(|\hat h_j|>K^2\).
In each remaining run, start a new interval whenever
the next step would bring the sum of absolute stepsizes above \(M=2K^2\).
Write \(\mathcal I\) for these intervals. We have
\begin{equation}\label{eq:count-partition}
 \#\{j:|\hat h_j|>K^2\}+|\mathcal I|\le Cn/K,\qquad
 V_I:=\sum_{j\in I}(1+|\hat h_j|)^{4/9},\qquad \sum_I V_I\le Cn .
\end{equation}
In fact \eqref{eq:count-quadratic} at \(4r\) gives
\(\#\{j:|\hat h_j|\ge r\}\le Cn/\sqrt r\) for \(1\le r\le K^2\).
Integration up to \(K^2\) bounds the retained absolute mass by \(CnK\)
and the sum of the \(4/9\)-moments by \(Cn\), since \(4/9<1/2\).
Each run contributes at most one final interval, and every other
interval has absolute mass greater than \(K^2\). This gives
\eqref{eq:count-partition}, since \(n\ge\kappa^q\ge K\).

\paragraph{Counting the steps above a threshold.}
Fix \(s=K/4\) and set \(t_j=\min(\hat h_j,s)\), so that
\(S_K^+(H)=4\sum_j(t_j)_+\). For \(t>0\) put
\(g(t)=t\min\{1/4,1024^2/\log^2(e+t)\}\) and for \(t\le0\) put
\(g(t)=0\). At a checkpoint we reserve \(g(t_j)\) of the step to
start the next component. We therefore compare the remaining sizes
\(v_j=t_j-g(t_j)\) with the threshold. For each interval \(I\) let
\[
 W_I=\sup_{z>0} z^q\#\{j\in I:v_j\ge z\}.
\]
Indeed \(v_j\ge3t_j/4\) for \(t_j>0\), so integrating this count gives
\begin{equation}\label{eq:count-moment}
 \sum_{j\in I}(t_j)_+
 \le \frac{(4/3)^q}{1-q}\,K^{1-q} W_I .
\end{equation}
The removed steps contribute only \(O(n)\) to \(S_K^+(H)\), since there
are \(O(n/K)\) of them and each contributes at most \(K\).
It remains to prove \(\sum_I W_I=O(n)\).

\paragraph{Estimating the count on one interval.}
Fix \(I\). A checkpoint set \(T\subseteq I\) partitions its unselected
indices into gaps, including the initial and terminal gaps. In gap \(U\)
write \(m_U\) for the largest signed prefix sum of the \(t_j\), including
the empty prefix, and \(d_U=m_U-\sum_{j\in U}t_j\) for the drop from
this maximum to the final sum. Let \(S_U=\sum_{j\in U}(t_j)_+\).
Write \(b_U=g(t_j)\) if checkpoint \(j\) follows \(U\), and \(b_U=s\)
for the terminal gap.

Let \(z_\star\) be a height attaining \(W_I\). When \(z_\star\) is
bounded, \(W_I=O(|I|)\), so only large \(z_\star\) matter.
As in~\cite[(5.6)--(5.9)]{JH}, take \(u\) between a sufficiently
large absolute constant and \(\log z_\star-4\sqrt{\log z_\star}\).
For the final comparison just below \(z_\star\), also take \(u=\log z_\star\)
as in~\cite[Lemma~5.2]{JH}. In either case put \(B=e^{u/32}\). For a gap \(U\) set
\(Y_U=B+m_U+b_U\) and \(Y_U^+=B+S_U+b_U\).
The first construction, which we call the ordinary test, uses one planar trajectory per checkpoint
(Appendix~\ref{app:tests}). Its lower bound on the logarithm of the
final-to-initial norm ratio, before the error estimated below, is
\(-A(T)\), where
\[
 A(T)=\log\frac{\prod_{U\ {\rm gap\ of}\ T}Y_U}
                   {s\prod_{j\in T}(v_j-B)}+\frac{C_p}{u}(|T|+1).
\]
Here \(C_p\) is a sufficiently large absolute constant. The sets \(T\)
are \(\{j:v_j>B+e^y\}\), their closed versions with \(\ge\),
and the sets obtained by adding checkpoints in~\cite[(5.24), (7.12)]{JH}.
At the top height we take \(T=\{j:v_j\ge z_\star\}\).
Every lower threshold selects every clipped step as well, since all
clipped steps have the same, largest value of \(v_j\). Hence the gap
steps use the original \(\hat h_j\).

For \(e^{u/2}\le B+e^y\le e^{u+\sqrt u}\), Jeong and Hong also use a
second construction. Recursively split each interval at the leftmost
index maximizing \(v_j>0\), leaving its earlier and later indices in
the two children. Stop when no positive entry remains. On certain gaps
the next \(\ell+2\) levels of splits are close to the equality case of
their scalar inequality, where \(\ell=\lfloor(\log_2u)/100\rfloor\).
On these gaps they use four gradient directions to improve the ordinary
lower bound. In our notation, the improvement in log norm is
\(\tfrac12\sum_{r\in\mathcal A(y)}(G_\ell(r)-C_p\omega_r)_+\).
Here \(\mathcal A(y)\) is the set of gap roots satisfying the adjacency
and near-equality conditions in~\cite[(6.45)--(6.49)]{JH};
\(G_\ell(r)\) is their weighted sum of changes in log ratio over the
subtree~\cite[(6.30)--(6.31)]{JH}, and \(\omega_r\) is the error in
this approximation~\cite[(6.53)]{JH}. We use their definitions with
heights \(v_j\) and gap sizes \(Y_U\). The factor \(1/2\) converts their
log squared-norm estimate to a log norm estimate.

Choose \(q<d<4/5\). Subtract \(e_I(y)=CV_Ie^{-dy}\) from each
test's log gain at its reference level \(y\). Let \(Q_I\) be the
larger of zero and the supremum of these adjusted gains. At bounded
\(z_\star\), set \(Q_I=0\).
We will prove
\begin{equation}\label{eq:count-scalar}
 W_I\le C\{V_I+K^q(1+Q_I)\}.
\end{equation}
The other ingredient is that the tests can be chosen independently
on different intervals and realized on one strongly convex objective.
By \eqref{eq:physical-simultaneous}, its log norm gain is at least
the sum of the positive adjusted gains up to an error of \(Cn/K^q\).
If \(D(H)\le1\), this log norm ratio is at most zero. If \(E(H)\le1\),
strong convexity and smoothness bound it by \(\tfrac12\log\kappa\).
Taking suprema gives
\begin{equation}\label{eq:count-aggregate}
 \sum_I Q_I\le Cn/K^q,
\end{equation}
since \(n\ge\kappa^q\).

\paragraph{The scalar estimate.}
We check how the boundary charges \(b_U\) and the adjustment \(e_I(y)\)
enter the proof of~\cite[Proposition~7.3]{JH}. Fix \(I\), write \(W=W_I\),
and assume \(W\ge |I|+K^q\), since otherwise \eqref{eq:count-scalar}
is immediate.

\emph{Splits.} At a split \(j\) the left and right children \(L,R\) receive the terms
\(g(t_j)\) and \(b_U\). With \(\tau=v_j-B>0\) the prefix identity
\(m_U=\max(m_L,\sum_Lt_i+t_j+m_R)\) gives
\[
 Y_L+Y_R+\tau-Y_U=\min(d_L,t_j+m_R),\qquad
 Y_L^++Y_R^++\tau-Y_U^+=0.
\]
Also \(Y_L\le Y_U\), since both \(g(t)\) and \(t-g(t)\) increase for
\(t>0\), and the checkpoint following \(U\) is either an ancestor
of \(j\) or is replaced by the terminal boundary term \(s\).
These identities verify the hypotheses of~\cite[Lemma~3.1]{JH} as in
their Lemma~3.3. At a split close to equality, in the precise sense of
\cite[Lemma~5.5]{JH} applied to both \(Y\) and \(Y^+\),
\cite[(5.15)--(5.16)]{JH} give
\(d_L=O(\delta^2\tau)\) and \(Y_L^+-Y_L=O(\delta\tau)\), where
\(\delta=u^{-1/10}\). The negative mass on the left is
\(Y_L^+-Y_L+d_L\); the maximal prefix passes the splitting step, so
\(d_R=d_U\). Replacing \(\tau\) by \(\tau-d_L\) gives
\(Y_U=Y_L+Y_R+\tau-d_L\). The corrected change in log ratio is
\(G_j=\log((\tau-d_L)Y_U/(Y_LY_R))\), as in~\cite[(5.27)--(5.31), (6.9)]{JH}.

\emph{Threshold integration.} Add \(Q_I+e_I(y)\) to the log potential
\(A(T)\) and its positive-mass counterpart, also for enlarged sets at
the same reference level. They are now nonnegative, since \(Y_U^+\ge
Y_U\), and the shifted potential bounds the extra gain of the second
test. The arithmetic--geometric mean estimate~\cite[(5.1)]{JH} for
\(\sum_U Y_U^q\) thereby incurs an endpoint cost at most
\(q e^{qb}(Q_I+e_I(b))\) at level \(b\). The jumps at a fixed level
are unchanged. In the integration by parts of~\cite[(5.4)]{JH} the
additional derivative cost on \([a,b]\) is at most
\[
 C\int_a^b e^{qy}|e_I'(y)|\,dy\le C V_I e^{-(d-q)a}.
\]
The same-level shifts cancel in the comparisons (5.24) and (7.12).
For the four-direction test, (6.66) bounds the gain pointwise by the shifted
potential, as needed in (6.70)--(6.71). Since \(a\ge u/4-O(1)\), the
terms involving \(e_I\), even multiplied by the polynomial window
factors, are summable over dyadic \(u\). The constant shift costs at
most \(qK^q Q_I\) at the endpoint and has zero derivative.

\emph{The terms reserved at checkpoints.} From \(\#\{t_j\ge z\}\le C Wz^{-q}\) one
obtains by integration
\[
 \sum_{t_j\ge e^U}g(t_j)^q\le CW U^{1-2q}.
\]
At an interior threshold \(z=B+e^y\) there are \(O(Wz^{-q})\) gaps
of total unselected positive mass \(O(Wz^{1-q})\). Concavity bounds
the sum of the \(q\)-powers of those masses by \(CW\); adding the
gap offsets, including the single terminal \(s\), preserves this
order. At a near-equality split in~\cite[(5.21)]{JH}, the child's
positive mass including its boundary term is comparable to \(\tau\).
If the boundary term supplies at least half of this
mass, charge \(\tau^q\) to that boundary instead of its steps.
The terminal charge is too large for such an interior child. At a
fixed boundary there is one left child and possibly a nested sequence
of right children. Their \(Y^+\) decrease under inclusion, and drop
by a factor \(1/(1+\sqrt2)+O(\delta)\) at each near-equality split. The sum
charged to boundary \(j\) is therefore \(O(g(t_j)^q)\). Here
\(g(t_j)\gtrsim e^{u/4-O(1)}\), so the total cost \(O(Wu^{1-2q})\)
is within (5.20). All remaining prefix comparisons in (6.61) use
\(Y_L\le Y_U\).

With these changes the right sides of the summed estimates
(5.11)--(5.13), (6.68)--(6.71), and (7.12) in~\cite{JH}
are controlled by \(W+K^qQ_I+V_I\) in place of \(W\).
Write \(\Phi(u)\) for the average of \(\sum_U(1+m_U+b_U)^q\)
at thresholds \(\exp(u+t\sqrt u)\), \(0\le t\le1\).
The weighted sum in \cite[Proposition~7.3]{JH} gives
\(\Phi(u)-\Phi(u/2)\le C(W+K^qQ_I+V_I)u^{-1/100}\). Changing the offset from
\(B\) to one there costs at most \(O(WB^q e^{-qu/2})\).
The top comparison at the closed set, as in \cite[Lemma~5.2]{JH},
gives \(\Phi(\log z_\star-4\sqrt{\log z_\star})\ge cW-C K^q Q_I-CV_I\).
At the bottom level \(\bar u\in[U,2U]\), the unselected positive
steps are bounded in terms of \(U\). Subadditivity and the boundary
charge estimate give
\(\Phi(\bar u)\le C_U|I|+K^q+CW U^{1-2q}\).
Summing the level differences and fixing \(U\) large absorbs the terms
proportional to \(W\), since
\(U^{1-2q}+\sum_{j\ge0}(2^j U)^{-1/100}\to0\).
This proves \eqref{eq:count-scalar}.

Summing \eqref{eq:count-scalar} over the intervals now gives
\(\sum_I W_I=O(n)\) by \eqref{eq:count-partition} and
\eqref{eq:count-aggregate}, since \(K^q|\mathcal I|\le Cn\).
The capped stepsize bound follows from \eqref{eq:count-moment}
and the contribution of the discarded steps.

\appendix
\section{Realizing the prescribed trajectories}\label{app:geometry}
We use the Moreau envelope~\cite{Moreau} to turn a convex function $P$ into a smooth objective with prescribed gradients. For a finite nonnegative convex function $P$ with $P(0)=0$, write $\operatorname{env}_1P(x)=\min_z\{P(z)+\|x-z\|^2/2\}$. For $0<\theta<1$, set
\[
 F_P(x)=\tfrac{\theta}{2}\|x\|^2+(1-\theta)\operatorname{env}_1P(x).
\]
We call $P$ the \emph{residual} in this representation. The resulting objective $F_P$ is $1$-smooth and $\theta$-strongly convex, with minimizer zero. To prescribe a gradient $g$ at a query $x$, it suffices to exhibit a supporting plane of $P$ at $z$ with slope $l$, where
\begin{equation}\label{eq:prox-interface}
 z=\frac{x-g}{1-\theta},\qquad l=x-z=\frac{g-\theta x}{1-\theta}.
\end{equation}
Indeed $l\in\partial P(z)$ says that $z$ is the minimizer defining $\operatorname{env}_1P(x)$, the \emph{proximal point} at $x$. As in Jeong and Hong~\cite[Lemma~4.1]{JH}, taking the maximum of $P$ and another convex function preserves the prescribed gradient whenever the other function is at most $P(z)$ at this point. This gives a way to check that GD follows two successive trajectories on one objective.

\subsection{Following a signed trajectory across checkpoints}\label{app:arc}

\paragraph{The planar trajectory.}
Ye and Liu~\cite[Section~3.1]{YLsmooth} construct a trajectory backwards by
projecting onto the convex hull of a circular arc and the origin. For signed
steps we use the disk of Jeong and Hong~\cite[Section~4.2]{JH}. For \(0<\sigma\le1/16\),
its center is \(z_c=(\sigma,-\sigma)\) and its boundary passes through
\(k=(\sigma^2,1)\). Intersect the disk with \(0\le\sigma^2 v_2\le v_1\)
and denote the resulting set by \(\mathcal K\). Its right endpoint on the
horizontal axis is \((c,0)\), \(1\le c\le1+2\sigma\). Let \(a=1-\theta\),
where \(0<\theta e^{4/\sigma}\le\sigma/4\). Write \(\Pi_{\mathcal K}\) for
Euclidean projection and \(\sigma_{\mathcal K}(z)=\max_{v\in\mathcal K}
\langle v,z\rangle\) for the support function. The gradient of its
Moreau envelope with parameter one is \(\Pi_{\mathcal K}\)~\cite[(4.12)]{JH}.
Adding the strongly convex quadratic gives the field
\(V(Q)=\theta Q+a\Pi_{\mathcal K}(Q)\).

We solve \(Q_{i-1}-h_i V(Q_{i-1})=Q_i\) backwards from \(Q_N=k\) for
any real steps with \(1-\theta h_i>0\). For this, take the region
between the positive horizontal axis and the following curve. Start
with the segment \([0,k]\), then follow \(\dot Q=V(Q)\) until its
projection first reaches \((c,0)\), say at \(Q_*=(X_*,W_*)\). Continue
with the ray
\(y=W_*(\theta x+ac)/(\theta X_*+ac)\).
To verify the construction, write \(Q=z_c+s(\cos\alpha,\sin\alpha)\)
along the circular part and let \(r\) be the disk radius. Then
\(V=z_c+(\theta s+ar)(\cos\alpha,\sin\alpha)\) and
\(s\dot\alpha=-\sigma(\sin\alpha+\cos\alpha)\). While
\(s\le e^{4/\sigma}\), direct substitution gives
\(0<\dot s\le1+4\sigma\) and \(-s\dot\alpha\ge\sigma(1-2\sigma)\).
The angular interval has length at most \(\pi/2+\sigma\). Integrating
\(d\log s/d\alpha\) gives \(X_*+W_*\le C e^{3/\sigma}\) before
the radius reaches \(e^{4/\sigma}\). The tangent slope decreases because,
with radial and tangential unit vectors \(e_\alpha,e_\alpha^\perp\),
\(\det(V,D\Pi_{\mathcal K}V)=(r/s)(V\cdot e_\alpha)
(V\cdot e_\alpha^\perp)<0\). Beyond \(Q_*\) the entire region
lies in the normal cone translated to \((c,0)\), as the ray's slope
is below \(W_* /(X_*-c)\). The boundary is therefore concave with
matching tangents, and the gradient is tangent on every face. The
backward equation now has a solution by the argument of Jeong and
Hong~\cite[Lemma~4.3]{JH}, who give a variational proof of the
signed-step construction of He and Cheng~\cite{HC}. Indeed, minimizing
\(\|z-Q_i\|^2/2-h_iG(z)\) over the region, where \(\nabla G=V\),
gives \(z-h_iV(z)=Q_i\) by tangency. A minimizer exists because
\(1-\theta h_i>0\) and the envelope term grows at most linearly.

Write \(Q_0=(X,W)\), \(\chi_j=\prod_{i\le j}(1-\theta h_i)\) and
\(m_\theta=\max_{0\le j\le N}(\chi_j^{-1}-1)/\theta\). Up to the
first \(j\) with \((Q_j)_1\le X_*\), the two coordinates satisfy
\((Q_j)_1+ac/\theta=\chi_j(X+ac/\theta)\), \((Q_j)_2=\chi_j W\).
Hence
\begin{equation}\label{eq:arc-prefix}
 X\le X_*+(\theta X_*+ac)m_\theta,\qquad
 W\le W_*(1+\theta m_\theta).
\end{equation}
The estimates hold on every suffix. If the positive mass is \(S\)
with \(\theta S\ll1\), then \(m_\theta\le S/(1-\theta S)\). If
\(U=\sum|h_i|\) and \(\theta U\ll1\), expanding the prefix products
instead gives \(m_\theta=m+O(\theta U^2)\), where \(m\) is the
largest signed prefix sum, including zero.

We also need a lower bound on \(X\) at one specified long step.
Whenever \(ac h_j\ge X_*\), take at that step the preimage
\(((Q_j)_1+ac h_j,(Q_j)_2)/(1-\theta h_j)\). Its first coordinate
is at least \(X_*\). It lies under the terminal ray because the whole
region lies below the affine extension of that ray. Telescoping the
first coordinate back to \(Q_0\), using \(\Pi_{\mathcal K}(Q)_1\in[0,c]\)
and \eqref{eq:arc-prefix} on all suffixes, gives
\(X\ge ac(h_j-\sum_{i<j}(-h_i)_+)-O(\theta U^2)\) when
\(X_*\le U,\ \theta U\ll1\).

\paragraph{Transfer at an ordinary checkpoint.}
We construct a component using the gradient field \(V\). Until we rescale
the objective at the end of this paragraph, \(h_i\) and \(b\) denote the
stepsizes for this field, with \(b\) the checkpoint step. Take
\(K_0\ge C(1+X_*+W_*/\sigma^2)\), \(\theta(K_0+S)\) sufficiently
small, and a checkpoint \(b>K_0\). Put
\(D=K_0+(\theta X_*+ac)S/(1-\theta S)\), \(Z=(1-X/D,W/D)\), and
\(J(x,y)=(x,-y)\). The points \(x_j=Z+JQ_j/D\) start at \((1,0)\)
and take the prescribed steps with gradients \(g_j=JV(Q_j)/D\).
Since \(V(k)=k\), the checkpoint sends the second coordinate to \((W+b-1)/D\).
To realize these gradients, take the convex function inside the envelope
in \eqref{eq:prox-interface} to be
\[
 P(z)=D^{-1}\max_{v\in\mathcal K}\langle Jv,z-Z/a\rangle
                    -(\theta/a)\langle Z,z\rangle .
\]
At the proximal point \(z_j=(x_j-g_j)/a\), the projection identity gives
the required subgradient \((g_j-\theta x_j)/a\). Equation~\eqref{eq:arc-prefix} on each suffix gives
\(D\ge X+W/\sigma^2+c+1\), \(P(0)=0\), \(\|Z\|+\|z_j\|\le C\),
and \((z_j)_1\ge0\). To see that \((z_j)_2\le W/(aD)\), observe that
reducing the second coordinate of any point of \(\mathcal K\) toward zero
preserves membership, so \((\Pi_{\mathcal K}Q)_2\le Q_2\).

To allow a positive input threshold, add a scalar quadratic of curvature
\(\alpha=[16(K_0+S)]^{-1}\) with initial value \(\alpha v_0^2=C_2\theta\),
\(v_0>0\). Its pre-checkpoint values \(v_j=v_0\prod_{t\le j}(1-\alpha h_t)\)
all exceed \(15v_0/16\). Its function inside the same envelope is
\(P_s(w)=(\alpha-\theta)w^2/[2(1-\alpha)]\), with proximal coordinate
\(w_j=(1-\alpha)v_j/a\). Set \(J_0=\sqrt{1+v_0^2}\),
\(e=(1,0,v_0)/J_0\) and \(o=(0,1,0)\), the input and output directions.
Let \(L\) be the maximum of the supporting affine functions of \(P+P_s\)
with the prescribed slopes at \(Y_j=(z_j,w_j)\).
All intercepts and slopes along \(o\) are nonpositive. Slopes along
\(e\) are at least \([-\theta Z_1+v_0(\alpha-\theta)v_j]/(aJ_0)
\ge2\theta/a\). Moreover \(L(Y_j)\ge c\alpha v_j^2-C\theta\),
enough to cover \(C\theta(1+v_0v_j)\) for the absolute constants
below. Choose \(C_2\) large and then \(\theta(K_0+S)\) small enough
that \(J_0\le2\).

For \(0\le d\le1\) translate the physical input by \(de\) and set
\[
 L_d(Y)=L(Y-de/a)-(\theta d/a)\langle e,Y\rangle+
                  \theta d^2/a-L(-\theta de/a).
\]
Then \(L_d(de)=0\), and the preceding margin makes \(L_d\ge0\)
at every translated proximal query. In \(L_d\) replace the \(e\)
coordinate by its maximum with \(d\), and the \(o\) coordinate by
its minimum with \(W/(aD)\). Take the maximum of the result with
zero to obtain \(U_d\). The slope signs make this a convex majorant
that preserves the supports at the queries. On \(te\) for \(t\le d\)
it is zero, and above the output threshold \(W/(aD)\) it is constant
in direction \(o\).

If \(A\) is the incoming coordinate value and \(s\le(A-s)/2\) its threshold, take
\(\ell=(A-s)/J_0,\ d=s/\ell\). Identify \(e\) with the preceding
output and place the perpendicular coordinates in a new orthogonal subspace.
Use the residual \(\ell^2 U_d(\cdot/\ell)\) for this gap. Its output
value and threshold are \(\ell(W+b-1)/D\) and \(\ell W/(aD)\).
At each query later residuals thus have zero subgradient on their
input ray. During the next gap the prior residual is flat on the
shared proximal coordinate. Minimizing over its private coordinates
leaves the prescribed active proximal point unchanged. By induction,
the sum of the gap residuals under one envelope realizes all the
transfers. The exact excess ratio is
\((b-1-\theta W/a)/(D J_0)\), and the outgoing threshold-to-excess
ratio is \(W/[a(b-1)-\theta W]\).

To obtain \eqref{eq:signed-block-transfer}, multiply this local objective by
\(1/4\). A step \(h_t\) on the scaled objective is a step \(h_t/4\) on the
local one, where we take \(\theta=16\mu\) and \(K_0=K/4\). With \(\sigma=32/\log K\) the
required bound on \(K_0\) follows from \(X_*+W_*\le C e^{3/\sigma}\).
Now \(\theta W/a\le1\), \(J_0\le e^{C\theta(K_0+S)}\), and
\(\theta X_*+ac\le1+C/\log K\). Substitution gives the asserted
transfer with no fixed cost per checkpoint. The threshold ratio is
\(O(W_*/b)=o(1)\) uniformly for the checkpoints used here, so
\(s/(A-s)\le1/2\) and \(J_0\le2\) give \(d\le1\) at the next gap.

\paragraph{Three quadratic coordinates distinguish the final query.}\label{app:spectral}
For \(r\) steps write \(P_j(\lambda)=\prod_{t=1}^j(1-\lambda h_t)\) for
\(0\le j\le r\), with \(P_0=1\). Use three quadratic curvatures
\(\lambda_k=\lambda_0+k\gamma\), \(k=-1,0,1\), spaced by \(\gamma>0\) in
\([4\mu,20\mu]\). Suppose none of the factors vanishes between the two
outer curvatures, endpoints included. A coordinate of the form
\(u_{j,k}=\xi_k+c_k P_j(\lambda_k)\) has final variable part
\(c_kP_r(\lambda_k)\). Choose its three magnitudes all equal to \(T\).
For a previous query \(j\), put \(p=P_r/P_j\). The magnitude of
the difference in coordinate \(k\) is exactly
\(T|1/p(\lambda_k)-1|\), regardless of the offsets \(\xi_k\).

Here is why three samples suffice even if signs alternate. Put
\(v_t=\gamma h_t/(1-\lambda_0h_t)\). Since \(1-\lambda h_t\) has no
zero on this interval, \(|v_t|<1\). Factoring the outer two multipliers gives
\[
 \log\frac{|p(\lambda_0)|}
               {\sqrt{|p(\lambda_{-1})p(\lambda_1)|}}
 =-\tfrac12\sum_{t>j}\log(1-v_t^2)\ge\tfrac12\sum_{t>j}v_t^2 .
\]
Within distance \(d<T/4\) of the last query, each of the three
\(\log|p(\lambda_k)|\) has magnitude \(O(d/T)\). Since
\(|v_t|\ge c(\gamma/\mu)\min\{|\mu h_t|,1\}\), this gives
\(\sum_{t>j}\min\{(\mu h_t)^2,1\}\le C(\mu/\gamma)^2d/T\).
In particular, for \(0<\eta\le1\) choose
\begin{equation}\label{eq:three-radius}
 d=cT\gamma^2\eta^2/(1+r).
\end{equation}
Then the sum is less than one, every term is untruncated, and
Cauchy--Schwarz gives \(\sum_{t>j}|h_t|\le\eta\).

For a marked interval use \(12\mu,16\mu,20\mu\). Positive steps
are at most \(\delta/\mu\) or exceed \(1/(4\mu)\), so every product
is nonzero throughout this interval. At each curvature the large
steps give at least a factor two, a small positive step costs at most
\(40\mu h_t\) in log distance, and negative steps cost nothing.
Start three independent quadratics at \(c_k=T/|P_r(\lambda_k)|\)
with unit initial norm. Then \(T\ge\min_k|P_r(\lambda_k)|/\sqrt3\).

To continue the trajectory after its last checkpoint, use the same curvatures.
In the original stepsize units the objective is \(\theta\|\cdot\|^2/2+
(L_0-\theta)\operatorname{env}_1 U\), where \(\theta=4\mu\),
\(L_0=1/4\). Suppose the output has just reached \(s+D\) with \(0\le s\le D\)
and its values at earlier iterates and their proximal points are at most \(s\).
Assume \(U\) is nonincreasing along this axis
and flat above \(s\) at every choice of the other coordinates.
These properties hold for the function \(U_d\) constructed above.
Section~\ref{app:tests} will also use the case where the other coordinates
are zero at the checkpoint and \(U\) vanishes on the positive output ray.
Identify the output axis with the span of a unit vector \(e\) with positive
entries in three dimensions. Add to \(U\) the convex function
\[
 \sum_{k=-1}^1 \frac{\lambda_k-\theta}{2(L_0-\lambda_k)}
                     (w_k-se_k)_+^2 .
\]
All three terms vanish at the earlier proximal points. During the waiting
interval of positive mass at most \(\delta/\mu\), write
\(t_k=\theta/\lambda_k\). The three new coordinates evolve as
\[
 u_{j,k}=(1-t_k)s e_k+(D+t_ks)e_k P_j(\lambda_k).
\]
Indeed each prefix product is at least \(1-20\delta\), so
\(u_{j,k}\ge(s+cD)e_k\). The proximal point of each hinge lies above
\(se_k\), where the preceding residual attains its minimum along the
output axis there. Its private coordinates can be minimized separately,
so each new gradient is \(\lambda_k u_{j,k}-(\lambda_k-\theta)se_k\).
In the simpler case, keep the private coordinates at zero. The
hinge proximal point projects onto the positive output ray, where
\(0\in\partial U\), so it solves the joint proximal problem as well.
Choose \(e_k=T/[(D+t_ks)P_r(\lambda_k)]\) with \(\|e\|=1\).
Then \(T\ge cD\), and \eqref{eq:three-radius} applies. Every query
before the checkpoint also lies at least \(cD\) from the final query,
since their projections onto \(e\) are at most \(s\) and at least \(s+cD\), respectively.
When the waiting interval ends the schedule, this final norm bound
alone suffices. An empty input chain starts with \(s=0,D=1\).

A variant handles the arbitrary strings between tested intervals in
Section~\ref{sec:proof-count}.
Let \(N\) count the steps \(h_t>4K^2\), in the original units.
Put \(\gamma=\mu/[16(1+r)]\). In the interval
\([4.5\mu,5.5\mu]\) exclude centers within \(2\gamma\) of any
positive reciprocal stepsize. These exclusions have total length at
most \(4r\gamma\le\mu/4\). For every remaining \(\lambda_0\) the
three curvatures lie in \([4\mu,6\mu]\) and no factor \(1-\lambda h_t\)
vanishes between them.
For every \(h>0\) the integral of \(\log^-|1-\lambda h|\) over
\([4\mu,6\mu]\) is \(O(\mu)\). Indeed, small \(\mu h\) gives a bounded
integrand, and otherwise substituting \(z=1-\lambda h\) gives
an upper bound \(2/h\). Averaging the sum over the three translates
and all \(h_t>4K^2\) therefore gives a center with
\(\sum_{k=-1}^1\sum_{h_t>4K^2}\log^-|1-\lambda_k h_t|\le C N\).
Since \(4K^2=o(1/\mu)\), each smaller positive step costs at most \(C\mu h_t\), while a
negative step costs nothing. Normalizing the initial coordinates
as in the marked interval gives
\(\log T\ge-C(1+N+\mu\sum_{0<h_t\le4K^2}h_t)\) and radius
\(cT\mu^2\eta^2/(1+r)^3\). For \(r=0\) take a unit initial point, since
there are no earlier queries to separate.

\paragraph{Starting the next portion on the same function.}\label{app:reset}\label{app:stages}
We give one construction for any finite list of trajectory portions.
For portion \(i\), let \(F_i\) be its \(4\mu\)-strongly convex,
\(1/4\)-smooth objective, minimized at zero and normalized by \(F_i(0)=0\).
Start the portion at a point of unit norm. At the end of every portion except
the last, a prescribed positive step \(4\le b_i\le1/(4\mu)\) is
reserved for the switch. Write \(u_i\) for the iterate just before it,
\(g_i=\|u_i\|>0\). Suppose that for every \(0<\eta\le\eta_0\)
any earlier query within \(d_i(\eta)\le g_i/8\) of \(u_i\)
starts a suffix whose stepsizes have absolute values summing to at most \(\eta\), excluding
\(b_i\). We will obtain in the next independent subspace the exact
prescribed starting point, scaled by \(A_{i+1}\), where
\[
 A_1=1,\qquad A_{i+1}=c_1b_i\mu^3 A_i d_i(\eta_i),\qquad
 \eta_i=c_2\min(\eta_0,b_i g_{i+1}).
\]
Use \(\eta_i=c_2\eta_0\) when the next portion is the last.
The loss will depend on each following \(g_{i+1}\) only once.

We first record the local modification used at each switch.
If \(K_0\) has the base curvature bounds and value and gradient
zero at zero, one can prescribe \(K_1=\lambda\|n\|^2/2\) on
\(\|n\|\le2r\) and \(K_1=K_0\) on \(\|n\|\ge R=C r/\mu\),
keeping curvature between \(2\mu\) and \(1/2\), for any
\(\lambda\in[4\mu,1/4]\). To verify this, subtract
\(\mu\|n\|^2\) from both specifications. Write \(H,p=\nabla H\)
for the resulting function values and gradients, with \(H_i,p_i\) their
values at \(x_i\), and let \(L=1/2-2\mu\). By the interpolation criterion
of Taylor, Hendrickx and Glineur~\cite[Corollary~1]{THG}, the required
inequality for each pair of prescribed values and gradients is
\(H_i-H_j-\langle p_j,x_i-x_j\rangle\ge\|p_i-p_j\|^2/(2L)\).
Only an outer \(x\) and an inner \(y\) need comparison. In either
order, the left side between \(x\) and zero is at least
\(\mu\|x\|^2/4+3\|p(x)\|^2/2\), by taking one quarter of its
strong convexity bound and three quarters of its smoothness bound.
Changing zero to \(y\) costs \(O(r\|x\|+r^2)\). For
\(\mu\le1/64\), \(\|p(x)-p(y)\|^2/(2L)\le
3\|p(x)\|^2/2+Cr^2\), and \(R=Cr/\mu\) absorbs the remainder.
Azagra and Mudarra~\cite[Theorem~1.3]{AM} show that the same pairwise
condition suffices on an arbitrary subset. It therefore gives an
\(L\)-smooth convex function with our prescribed values and gradients
on the two regions. Adding back \(\mu\|n\|^2\) gives \(K_1\).

Perform the following construction on each scaled base
\(F(u)=A_i^2 F_i(u/A_i)\) separately. Abbreviate
\(u_*=A_iu_i,\ d=A_i d_i(\eta_i),\ b=b_i\).
Take a copy of its space, initially zero, and rotate \((u,w)\)
to \((t,n)= (u+\epsilon w,w-\epsilon u)/s\), where
\(s=\sqrt{1+\epsilon^2}\) and
\(\epsilon\|u_*\|/s=r=c_0\mu^2d\). Define
\[
 T(t)=F(st)/s^2,\qquad K_0(n)=\epsilon^2 F(-sn/\epsilon)/s^2.
\]
The uncoupled iterates of \(T+K_0\) are
\((u_j/s,-\epsilon u_j/s)\) for the original iterates \(u_j\).
Write \(t_*=u_*/s,\ n_*=-\epsilon u_*/s\). Choose smooth cutoffs
\(\xi(t)\), \(\beta(n)\) that equal one within \(d/(4s)\) of \(t_*\)
and \(r/4\) of \(n_*\), respectively, and vanish beyond twice
these radii. Choose \(K_1\) as above with \(\lambda=1/b\).
Use the modified base
\(T(t)+K_0(n)+\xi(t)(K_1(n)-K_0(n))\).
For the next portion's full vector \(Z\) add the adjacent coupling
\[
 -\xi(t)\beta(n)\langle v,Z\rangle\,\zeta(Z/r),\qquad
 \|v\|\le 2 A_{i+1}/b,
\]
where \(\zeta\) is a fixed smooth cutoff equal to one on the unit
ball and vanishing outside the ball of radius two. The last portion
keeps its scaled base.

\(K_1-K_0\) and \(\nabla K_1-\nabla K_0\) vanish outside
\(R\) and are \(O(R^2)\) and \(O(R)\) inside. Hence the modified
base has an a.e.\ block-diagonal Hessian part with spectrum in
\([2\mu,1/2]\) plus an error of norm \(O(R/d+R^2/d^2)\le\mu/4\)
for small \(c_0\). The coupling has Hessian norm
\(O(\|v\|/r)\le\mu/8\) for small \(c_1\). On the direct sum, each
coordinate space is incident to at most two couplings. Thus their
total quadratic-form error is at most \((\mu/4)\sum_i\|p_i\|^2\)
on any vector \((p_i)_i\). The full objective has curvature between
\(\mu\) and one and is stationary at zero.

It remains to choose \(v\) so that the reserved step reaches the desired
point in the next subspace. Compare the modified trajectory to that of
\(T+K_0\), called the reference trajectory here. Until its first query within \(d\)
of \(u_*\) the modifications at this boundary vanish. From that query
to the reserved step the absolute stepsizes sum to at most \(\eta_i\).
The \(n\)-coordinate of the reference trajectory then has norm at most
\(2r\), since the reference states start within \(d\) of \(u_*\) and grow by at most
\(\exp(\eta_i/4)\). While the actual \(n\)-coordinate is within \(r\)
of its reference value, the gradient perturbation is \(O(r)\).
Indeed, \(K_0,K_1\) and their gradients are \(O(r^2)\) and \(O(r)\) there,
respectively. Iterating the Lipschitz gradient bound gives deviation
\(C\eta_i r e^{C\eta_i}<r/8\) for small \(\eta_0\), so the estimate holds
up to the reserved step. The new coordinates separately satisfy
\(\|Z\|\le C\eta_i\|v\|\). If the next portion has a reserved step,
its cutoff \(\xi\) vanishes
throughout \(\|Z\|\le c A_{i+1}g_{i+1}\). The choice of \(\eta_i\)
keeps \(Z\) there for every trial \(v\), so all later portions
stay at zero.

At the reserved query both cutoffs equal one, \(\|n\|<2r\), and
\(\zeta(Z/r)=1\). The gradient in the \(n\)-coordinates equals \(n/b\). The step
sets \(n=0\), where \(\beta\) vanishes on a neighborhood. It stays
zero thereafter. The next state satisfies \(Z^+(v)=bv+Z-b\nabla G(Z)\)
for its modified base \(G\), so \(\|Z^+(v)-bv\|\le
C\eta_i(1+b)\|v\|\le b\|v\|/4\). This is a continuous map of \(v\).
For the target \(V\) of norm \(A_{i+1}\), Brouwer's theorem applied to
\((V-Z^+(v)+bv)/b\) on \(\|v\|\le 2A_{i+1}/b\) gives \(Z^+(v)=V\).
Choose these vectors in chronological order. Later couplings were
inactive for every previous trial, and earlier couplings vanish for
the rest of the run. This establishes all the prescribed trajectories
on one fixed objective.

Let \(r_i\) count the steps of portion \(i\) before its reserved step.
Suppose \(g_i\ge cG_i\) and
\(d_i(\eta)\ge cG_i\mu^2\eta^2/(1+r_i)^3\) for each portion
ending in a switch, while the last portion ends at norm at least
\(cG_{\rm last}\) from unit input. These hypotheses include
\eqref{eq:three-radius} in all the above cases. Taking logarithms
of the amplitude recursion and using \(b_i\ge4\) yields
\begin{equation}\label{eq:joins-gain}
 \log\frac{\|x_n\|}{\|x_0\|}\ge
 \sum_i\log G_i-
 C\left[1+\sum_{i\ {\rm with\ switch}}\!
       \left(\log\kappa+\log(1+r_i)+\log^- G_{i+1}\right)\right].
\end{equation}
Here \(\log^-x=\max(0,-\log x)\). The term for the last portion can
be omitted from the negative-part sum. In particular the cost of a
switch is independent of the number of earlier or later switches.
\subsection{Constructing the trajectories for Section~\ref{sec:proof-count}}\label{app:tests}
Fix one interval of Section~\ref{sec:proof-count}. Write $h$ for a local stepsize, equal to one quarter of the original stepsize. We return to the original units when we join intervals. The sum of absolute local stepsizes is at most $M$. Use $F_P$ with $\theta=16\mu$ and write $a=1-\theta$. Then the objective $F_P/4$ used with the original steps is $1/4$-smooth and $4\mu$-strongly convex, as required in \eqref{eq:joins-gain}. As in Jeong and Hong~\cite[Sections~4.1, 6.5]{JH}, we construct the function backwards. A new piece ends at the starting point of the suffix already constructed. The identity \eqref{eq:prox-interface} lets us verify the gradients on both pieces by comparing supporting planes.

\paragraph{An ordinary checkpoint.}
Fix a threshold $B+Y$, with $B=e^{u/32}$ and $Y=e^y$ as in Section~\ref{sec:proof-count}. Let $b$ be the selected local stepsize, $t=\min(b,K/4)$ its capped value, and $m,U$ the largest signed prefix sum and the sum of absolute stepsizes in its preceding gap. The norm ratio prescribed by this checkpoint is
\[
 A_b=\frac{t-g(t)-B}{B+m+g(t)}.
\]
Every clipped step is a checkpoint, including for the enlarged sets in \cite[(5.24), (7.12)]{JH}, so the capped and original local steps agree throughout each gap. At every selected step $t-g(t)-B\ge Y/u^C$. For the second enlargement this follows from the range $k\le\ell+2$ used with (7.12), where $\ell=\lfloor(\log_2u)/100\rfloor$. Throughout, $u/4-O(1)\le y\le u+3\sqrt u$.

Use the planar orbit of Section~\ref{app:arc} with $\sigma=\min\{256/u,768/\log(e+t)\}$. Write $(X,W)$ for its starting point, $(\sigma^2,1)$ for its final query, and $c_\sigma=1+O(\sigma)$ for the positive endpoint of its gradient set on the first axis. Reflect the second coordinate, divide by $D$, and translate by $Z$, where
\[
 D=X+\sigma^2(b-1),\qquad Z=(1-X/D,W/D).
\]
Now the gap starts at $(1,0)$ and the step $b$ ends at $(0,R)$ with $R=(W+b-1)/D$. The signed prefix estimate of Section~\ref{app:arc} gives $X\le X_*+(1+C/u)m+C\theta M^2$ and $W\le(1+C\theta M)W_*$, where $X_*+W_*\le C e^{3/\sigma}$. The required $\theta e^{4/\sigma}\le\sigma/4$ follows from $t\le K/4$ and $u=O(\log K)$. Thus for $t\le e^{3u}$ we have $X_*\le B$ and $W\le C Y^{.05}$. Above that range $X_*+W_*\le C(e+t)^{1/256}$. In both cases $W\le\sigma^4(b-1)$ and $c_\sigma<\sigma^2(b-1)$ for large $u$, and $Y/u^{C'}\le D\le CM$.

Let $J(x,y)=(x,-y)$ denote the reflection and let $\mathcal K$ be the gradient set of Section~\ref{app:arc}. In the representation \eqref{eq:prox-interface}, the corresponding planar convex function is
\[
 P_0(z)=D^{-1}\max_{v\in\mathcal K}\langle Jv,z-Z/a\rangle
                 -(\theta/a)\langle Z,z\rangle .
\]
It satisfies $P_0(0)=0$ by $v_1\ge\sigma^2v_2$ and $W\le\sigma^4(b-1)$. At a planar query its proximal point has $z_1\ge0$, $z_2\le W/(aD)$, and $\|z\|\le CM$, by the projection formula of Section~\ref{app:arc}. The second coordinate of each supporting slope of $P_0$ at these queries is nonpositive, and $P_0(z)\ge-C\theta M$ by testing $v=0$.

We add one scalar coordinate $w$. Its contribution makes the supporting plane of the current piece active at each query in the gap. We also need the checkpoint to send $w$ to zero so that the next piece starts on its input axis. Set $\alpha=1/[16(1+b+U)]$ and let $T>0$ be the value of $w$ just before the checkpoint, to be chosen below. On the positive axis take a continuous derivative $k$ with $k(0)=0$, successive slopes $\alpha,1/2,\theta$, and breakpoints
\[
 v_*=T\frac{1/2-1/b}{1/2-\alpha},\qquad T.
\]
Set $k(-w)=-k(w)$. Since $k(T)=T/b$, the step $b$ sends $T$ to zero. Obtain the earlier iterates $w_j$ by inverting the gap steps from $T$. For $h>0$ the map $w\mapsto w-hk(w)$ has a positive preimage at least as large as its positive target, by continuity and $1-\theta h>0$. For $h<0$ the preimage is at least the smaller of $v_*$ and the target divided by $1+\alpha|h|$. Hence every $w_j$ is at least $v_*e^{-\alpha U}\ge T/3$. If $w_0>T$, stop at the first $w_j\le T$. Until that step all gradients are evaluated on the last linear segment, so $w_j+d=(w_0+d)\prod_{i\le j}(1-\theta h_i)$ with $d=T(1/b-\theta)/\theta$. Inverting this product gives
\begin{equation}\label{eq:physical-clock}
 w_0\le T\bigl(1+(m+C\theta M^2)/b\bigr).
\end{equation}
The same bound holds at each query with $m$ replaced by the sum of positive gap steps. Define $P_s$ by $P_s(0)=0$ and slope $(k(w)-\theta w)/a$ at $z_s=(w-k(w))/a$. The map $w\mapsto z_s$ is increasing onto $\mathbb R$ and the specified slopes are nondecreasing, so $P_s$ is convex. At the queries, $c\alpha T\le k(w)-\theta w\le T/b$. Evaluating on the first segment at $w=T/3$ also gives $P_s(z_s)\ge c\alpha T^2$, since $\theta\ll\alpha$.

The initial point is $J_1 e$, where $J_1=\sqrt{1+w_0^2}$ and $e=(1,0,w_0)/J_1$. Take $o=(0,1,0)$ as the input direction of the suffix, and scale the suffix to start at $Ro$. Suppose its convex function in \eqref{eq:prox-interface} is at most $k'Rz_+$ at $zo$, for some $k'\ge0$, and its proximal points have nonnegative projection on $o$. At the current proximal points its value is at most $Ck'bW/D^2$. Choose
\begin{equation}\label{eq:physical-amplitude}
 T^2=C\alpha^{-1}(\theta M^2+k'bW/D^2)
\end{equation}
and define $P$ as the maximum of the suffix function, zero, and the supporting planes of $P_0+P_s$ at the current proximal points. At these points $P_s$ covers both the negative part of $P_0$ and the value of the suffix function. Conversely, the new planes have nonpositive intercept and nonpositive slope along $o$, so lie below the suffix function at its proximal points. Both trajectories therefore keep their gradients, and the checkpoint ends at $Ro$. The proximal points have nonnegative projection on $e$, and along this input direction
\begin{equation}\label{eq:physical-input}
 P(ze)\le C J_1 z_+(D^{-1}+b^{-1}+\theta).
\end{equation}
Indeed, the planar slope on the first coordinate lies between $-\theta Z_1/a$ and $c_\sigma/D+\theta\|Z\|/a$. The scalar contribution makes the total input slope nonnegative since $w_0(k(w_j)-\theta w_j)\ge c\alpha T^2$, and $w_0T/J_1^2\le3$. All later pieces are perpendicular to $e$.

For interior joins take a common $k'=u^{C'}/Y$ large enough for \eqref{eq:physical-input} and the four-direction input bound below. Equations \eqref{eq:physical-clock}--\eqref{eq:physical-amplitude} give
\[
 DJ_1\le D+C\sqrt\theta\, M^{7/2}
       +C(b+m+C\theta M^2)\sqrt{(1+b+U)Wk'/b}.
\]
When $t\le e^{3u}$ and $U\le Y^{9/5}$ the square root is $O(u^{C''}Y^{-.075})$. Multiply by $t/b$, use $g(t)\gtrsim t/u^2$, and recall the bound on $X$. For large $u$ we have $\sigma^2t\le(3/4)g(t)$ by the choices of $g$ and $\sigma$. The remaining quarter of $g(t)$ absorbs the term proportional to $t$ in the error. This gives $(t/b)DJ_1\le (B+m+g(t))e^{C/u}$. Since $(t/b)(W+b-1)\ge t-g(t)-B$, the actual gain $R/J_1$ is at least $A_b e^{-C/u}$. Above either cutoff the same inequality loses at most an extra factor polynomial in $1+U$ or $e+t$, respectively, by the displayed bound for $DJ_1$ and $X_*+W_*\le C(e+t)^{1/256}$ when $t>e^{3u}$. The factor $t/b$ permits $b$ to exceed its capped value $t$ by an arbitrary factor.

\paragraph{Preparing the next gap for four gradient directions.}
Put $\rho=1+\sqrt2$ and $\delta=u^{-1/10}$. On a gap indexed by $v\in\mathcal A(y)$, we use the four directions of Jeong and Hong~\cite[Lemmas~6.1--6.3]{JH}. We first show that the next gap satisfies the two conditions that permit this construction. In their notation~\cite[(6.33)]{JH}, if the suffix begins at $Y_0$, its convex function $P$ must satisfy, for constants $H>0$ and $Z_0\ge0$,
\begin{equation}\label{eq:physical-interface}
 P(sY_0)\le s_+\|Y_0\|^2/H,\qquad
 0\le\langle Y_0,w\rangle,\quad
 (\langle Y_0,w\rangle-Z_0\|Y_0\|^2)_+\le 8H P(w)
\end{equation}
at every suffix proximal query $w$, with the ray bound valid for all real $s$.

Let $r$ index the gap immediately following the checkpoint after $v$, as in~\cite[(6.48)]{JH}. Suppose we use the ordinary construction on that gap and write $j$ for its following checkpoint. Its positive steps sum to at most $CY$ by the near-equality estimate for the positive sums. Its drop $d_r$ is at most that of the left child of $j$, since the two intervals have the same right endpoint. Thus $d_r\le C\delta^2\rho^\ell Y$ by~\cite[(5.15), (6.47)]{JH}; the sum of absolute stepsizes satisfies $U=O(Y)$. The step at $r$ has size $(1+o(1))(v_r-B)\asymp Y$, and its preceding negative steps have total absolute size $O(\delta h_r)$ by \cite[(5.16)]{JH}. Section~\ref{app:arc} gives $X\ge(1-o(1))h_r$. Writing $Y_r=B+m_r+g(t_j)$ for $Y_U$ on the gap indexed by $r$, we obtain $D\asymp Y$, $D\le(1+o(1))Y_r$, $b\le C\rho^\ell Y$, and $Z_1\le C\rho^\ell/u^2$. These steps are uncapped by \cite[(5.6)--(5.9), (6.46)--(6.49)]{JH}. Equations \eqref{eq:physical-clock}--\eqref{eq:physical-amplitude} give $\max w_j\le CT$ and $T^2\le C\theta M^3+u^C Y^{-.15}$.

To supply the lower bound in \eqref{eq:physical-interface}, also include the hinge
\[
 (c_\sigma J_1/D)(\langle e,z\rangle-z_*)_+,\qquad
 z_*=(Z_1/a+w_0\max_j z_{s,j})/J_1 .
\]
At a present query its positive affine part is at most $(c_\sigma/D)(z_1-Z_1/a)$, which is supplied by using $(c_\sigma,0)$ in the definition of $P_0$ and bounding its negative linear term with $P_s$. Its value at every subsequent query is zero. The normalized input slopes of all planes, including the hinge, are at most $6/D$ by $D/b=O(1)$ and the scalar correction $O(T^2/b)$. Thus \eqref{eq:physical-interface} holds at $Y_0=J_1e$ with $H=D/6$ and $Z_0=z_*/J_1\le u^{-3/2}$; moreover $(v_r-B)/(8\rho^\ell)\le H\le Y_r/5$.

\paragraph{Four directions under quadratic contraction.}
Fix a gap whose root $r$ belongs to $\mathcal A(y)$. The $2^\ell$ intervals at depth $\ell$ of its tree are the cells. Each cell has four successive phases, each consisting of a gap and its checkpoint, as in~\cite[(6.12), (6.55)--(6.58)]{JH}. Their common scale is $M_c=\rho^{-\ell-2}Y_r$. Within one cell start $\chi_0=1$ and write $\chi_i=\prod_{t\le i}(1-\theta h_t)$, $\bar h_i=a h_i/\chi_i$, and $\chi$ for the final product. For phase $j=0,\ldots,3$, let $b$ be its checkpoint stepsize, $\bar b$ its rescaled value, and let $\bar m,\bar d$ be the maximum prefix sum and the drop to the final prefix sum for the steps $\bar h_i$ in its gap. Define
\[
 M_c\mu_j=B+\bar m+g(b),\qquad M_c k_j=\bar b-g(b)-B-\bar d .
\]
Thus $\mu_j+k_j$ is the sum of the rescaled stepsizes in phase $j$, divided by $M_c$. Each prefix sum $\xi$ of the rescaled gap steps satisfies $\xi+\chi_i^{-1}\le M_c\mu_j$ since $\chi_i^{-1}\le2\le B$. All the steps here are uncapped. Replacing the rescaled steps by the original local steps changes every normalized parameter by $O(\theta M^2/M_c)$.

Write $\tau_j=\mu_j+k_j$. First set $\theta=0$. For each gap indexed by $w$, let $j(w)$ be its following checkpoint and $q_w=(v_{j(w)}-B-d_w)/Y_w$. At a near-equality split $w$ with children $L,R$, the right child retains $d_w$ and its boundary term. This gives $q_Lq_R/q_w=e^{G_w}$, where $G_w$ is the corrected change in the logarithm at this split, as in~\cite[(6.9)--(6.10)]{JH}. The near-equality estimates \cite[(6.51)--(6.60)]{JH} therefore apply with the boundary terms $b_U$. They give
$\mu_j=1+O((\ell+2)\delta)$, $(k_0,k_1,k_2)=(\rho^{-1},1,\rho^{-1})+O((\ell+2)\delta)$ and $1-o(1)\le k_3\le C\rho^{2\ell}$, with the perturbation above absorbed. Let $N$ be the largest negative mass in the first three barred gaps. Then $N=O(\delta M_c)$ by~\cite[(5.16)]{JH}.

To verify \eqref{eq:physical-interface} between consecutive cells, use $(H,Z_0)=(M_c,0)$. The last cell uses the ordinary entrance described above, or $(\rho^{-\ell-2}Y_{\rm next},0)$ if the next gap also uses four directions, where $Y_{\rm next}$ is that gap's size. By the entrance bounds $cM_c\le H$ and $H/(M_ck_3)\le(\rho+O(\delta))/5<1/2$, as in \cite[(6.59)--(6.61)]{JH}. Indeed $H\le Y_{\rm next}/5$ also at an internal boundary since $M_c/Y_{\rm next}=(1+o(1))\rho^{-2}$. Passing left from the right child of the common checkpoint to the successor can only decrease its prefix maximum and charge. Put $\zeta=k_3 Z_0$ so that $\zeta=O(1/u)$, nonzero only in the last cell.

Take the four unit gradient directions $\varphi_j$ and the phase starting points $Q_j$ from \cite[Lemma 6.1]{JH}, with $e=Y_0/\|Y_0\|$, $\varphi_3=-e$ and $\langle e,\varphi_j\rangle=d_\varphi$ for $j<3$. In these parameters $d_\varphi\to1/(2\rho)<1/4$, so the preceding bounds give $2\chi k_3 N\le H$ and $8d_\varphi H\le M_c\chi(k_3-\zeta)$ with slack. Set
\[
 \lambda=\frac{\|Y_0\|}{\chi(k_3-\zeta)},\qquad
 \psi(z)=\frac{\lambda^2}{M_c}\big(\max_j\langle\varphi_j,z/\lambda\rangle-d_\varphi\chi\zeta\big)_+,
 \qquad \Psi=\max(P,\psi).
\]
Extend $P$ by orthogonal projection onto the suffix space and take all new directions perpendicular to the suffix except $e$. At a phase-$j$ prefix the intended point is $x=\lambda\chi_i(Q_j-\xi\varphi_j/M_c)$. Its intended proximal point $w=x-\lambda\varphi_j/M_c$ equals $\lambda\chi_i\widehat z$ with
$\widehat z=Q_j-\mu_j\varphi_j+s\varphi_j$, $s=\mu_j-(\xi+\chi_i^{-1})/M_c\ge0$.
By \cite[Lemma 6.1]{JH}, plane $j$ maximizes there. For $j<3$, comparison with $-e$ at $s=0$ gives $A=\langle\varphi_j,\widehat z\rangle\ge d_\varphi\mu_j$, hence $A-d_\varphi\chi\zeta/\chi_i\ge d_\varphi/2$. Also $\langle e,\widehat z\rangle_+\le d_\varphi N/M_c$, so the first inequality in \eqref{eq:physical-interface} and $2\chi k_3N\le H$ give $P(w)\le\lambda^2\chi_i d_\varphi/(2M_c)\le\psi(w)$. For $j=3$, $\widehat z=-(\zeta+s)e$, hence $P(w)=0$ and the active plane lies at or above the cutoff since $s+\zeta(1-d_\varphi\chi/\chi_i)\ge0$. Thus $\lambda\varphi_j/M_c\in\partial\Psi(w)$. Dividing the GD update by $\lambda\chi_{i+1}$ now gives the step $\bar h_{i+1}$, ending the fourth phase at $Y_0$. Conversely at a suffix query $\psi(w)=(\lambda d_\varphi/M_c)(\langle e,w\rangle-\lambda\chi\zeta)_+\le P(w)$ by $8d_\varphi H\le M_c\chi(k_3-\zeta)$ and $\lambda\chi\zeta\ge Z_0\|Y_0\|$. Both sets of iterates are therefore retained.

At the new input $\lambda Q_0\perp e$, \eqref{eq:physical-interface} holds with $(M_c,0)$. Indeed $Q_0=\sum_{j<3}\tau_j\varphi_j+(\mu_3+\zeta)\varphi_3$ has nonnegative inner products with all $\varphi_j$, bounded by $\|Q_0\|^2$ since $\tau_j>1$ and $Q_0\perp e$. At current queries of the first three phases,
$0\le\langle Q_0,\widehat z\rangle\le7A\le8(A-d_\varphi\chi\zeta/\chi_i)$. The first inequality follows by expanding $\widehat z=Q_3+(k_j+s)\varphi_j+\sum_{j<l<3}\tau_l\varphi_l$, the second from $\sum_{j<3}\tau_j+\mu_3+\zeta\le7$ and the maximizing-plane property. The last follows from $\mu_j\ge3/5$ and $\chi\zeta/\chi_i\le1/32$. The remaining proximal points are perpendicular to $Q_0$.

The squared norm multiplier per cell is $\chi^2(k_3-\zeta)^2/\|Q_0\|^2$. First evaluate the calculation of Jeong and Hong~\cite[(6.20), (6.28)--(6.32)]{JH} with the original local stepsizes and $\zeta=0$. Their formula for $\|Q_0\|$ is smooth and bounded away from zero near the equality case, and independent of $k_3$ for fixed $\zeta$. Thus the total change in log squared gain is $O(2^\ell\theta M^2+|\zeta|)=O(1/u)$. The errors in (6.29) add to the term in $\omega_r$ involving the descendants, while $d_r/(v_{j(r)}-B)$ converts $q_r$ to the ordinary factor. We obtain precisely the bound needed in Section~\ref{sec:proof-count},
\begin{equation}\label{eq:physical-enhanced}
 \log\frac{\|x_{\rm end}\|^2}{\|x_{\rm start}\|^2}
 \ge 2\log\frac{v_{j(r)}-B}{Y_r}+G_\ell(r)-C(\omega_r+1/u).
\end{equation}

\paragraph{One objective for all the intervals.}
On each interval $I$ choose any allowed test, independently of the other intervals. Write $S_I$ for its log gain after subtracting $e_I(y)$ as in Section~\ref{sec:proof-count}, and take $S_I=0$ on intervals without a test. For every such choice on a stable schedule we construct one admissible objective satisfying
\begin{equation}\label{eq:physical-simultaneous}
 \log\frac{\|x_n\|}{\|x_0\|}\ge\sum_I(S_I)_+-Cn/K^q .
\end{equation}
Write $J=|\mathcal I|$. The tests use cap $s=K/4$ and terminal charge $s$ in local units. The stepsizes $h_t$ below are in physical units.

First multiply the gains within one interval. Bound \eqref{eq:physical-enhanced} adds half the lower bound on the log squared-norm improvement for each gap using four directions. The $C/u$ costs fit the penalty $C_p(|T|+1)/u$ in $A(T)$, and the terminal factor $s/(s+B+m_{\rm tail})$ is at most one. For disjoint gaps of local absolute mass $U$, concavity gives $\sum U^{4/9}\le V_I$. Integrating $\#\{U\ge z\}\le V_I z^{-4/9}$ above $Y^{9/5}$, including the endpoint term, bounds the large-gap losses by $C(1+y)V_Ie^{-4y/5}$. The steps above $e^{3u}$ contribute at most $C(1+u)V_Ie^{-4u/3}$. Both are covered by $e_I(y)$ since $d<4/5$ and $y\le u+3\sqrt u$.

Discard intervals with $S_I\le0$. In each remaining interval reserve the first and last checkpoints of the chosen set to connect the interval to the preceding and following trajectories. Keep the middle pieces, making the last of them ordinary with empty suffix ($k'=0$). Include the additional hinge above if the preceding piece uses four directions. These changes cost $O(\log\kappa)$ per interval. An ordinary factor is at most $M$. A gap using four directions has log norm gain at most $C(1+\ell)2^\ell$ by the cell multiplier and $k_3\le C\rho^{2\ell}$; hence \eqref{eq:physical-enhanced} bounds its counted gain by this amount after subtracting the stated errors. Also $2^\ell\le u^{1/100}$, and a replacement ordinary factor is at least $\kappa^{-C}$ by its explicit $R,J_1$ and $k'=0$. Any interval with fewer than four pieces contributes only $O(\log\kappa)$ and may be skipped.

The last ordinary piece ends at $Ro$, where all coordinates perpendicular to $o$ are zero. Its planes have nonpositive intercepts and nonpositive slopes along $o$. The added hinge, if present, and the preceding pieces use directions perpendicular to $o$, so the full residual vanishes on the ray $\{zo:z\ge0\}$. Before this endpoint, the physical iterates and their proximal points project onto $o$ at coordinates at most $s_o=W/(aD)$. Moreover $s_o/(R-s_o)=W/[a(b-1)-\theta W]\le1$ because $W\le\sigma^4(b-1)$. Both properties persist under scaling.

Use the three quadratic coordinates from Section~\ref{app:arc} to continue the trajectory through the final open gap up to the last reserved checkpoint. The sum of its absolute stepsizes in the original units is at most $4M=o(\kappa)$, and every reserved step belongs to $[4,\kappa/4]$ for large $\kappa$. For unit initial norm, the lower bound $G_i$ on the final norm satisfies $G_i\ge\exp(S_I-C\log\kappa)$, since the output exceeds the threshold by at least half its norm. At the last query the radius in \eqref{eq:three-radius} is at least $cG_i\mu^2\eta^2/(1+r_i)$, where $r_i$ is the number of steps in this final gap.

Between retained intervals, as well as before the first and after the last, use the three quadratics of Appendix~\ref{app:arc}. On each such portion their total loss in log norm is at most $C(1+\#\{t:h_t>4K^2\}+\mu\sum_{0<h_t\le4K^2}h_t)$, and the corresponding radius is at least $cG_i\mu^2\eta^2/(1+r_i)^3$. An empty portion has no earlier query to separate. If no interval survives, use the three quadratics throughout.

To apply \eqref{eq:joins-gain}, sum the losses on the retained intervals and the quadratic portions. This sum also bounds $\sum_i\log^-G_i$, since each retained interval has $S_I>0$. By \eqref{eq:count-partition} there are $O(n/K)$ portions and discarded steps in total, and the retained absolute mass is $O(nK)$. The scanned lengths sum to at most $n$, so concavity gives $\sum_i\log(1+r_i)=O((n/K)\log(1+K))$. The total loss from $\sum_I(S_I)_+$ in \eqref{eq:joins-gain} is therefore $O((n/K)\log\kappa+\mu nK)=O(n/K^q)$, as claimed.
\par\smallskip
\noindent\emph{AI disclosure.} OpenAI Codex assisted in developing the proof and writing the manuscript.

\begin{thebibliography}{THG17}
\small
\raggedright
\setlength{\itemsep}{.25em}
\bibitem[AM17]{AM}
D. Azagra and C. Mudarra.
An Extension Theorem for Convex Functions of Class $C^{1,1}$ on Hilbert Spaces.
\emph{Journal of Mathematical Analysis and Applications} 446, 2017.
\href{https://doi.org/10.1016/j.jmaa.2016.09.015}{doi:10.1016/j.jmaa.2016.09.015}.
\bibitem[AP25a]{AP}
J.~M. Altschuler and P.~A. Parrilo.
Acceleration by Stepsize Hedging: Multi-Step Descent and the Silver Stepsize Schedule.
\emph{Journal of the ACM} 72(2), article 12, 2025.
\href{https://doi.org/10.1145/3708502}{doi:10.1145/3708502}.
\bibitem[AP25b]{APconv}
J.~M. Altschuler and P.~A. Parrilo.
Acceleration by Stepsize Hedging: Silver Stepsize Schedule for Smooth Convex Optimization.
\emph{Mathematical Programming} 213, 1105--1118, 2025.
\href{https://doi.org/10.1007/s10107-024-02164-2}{doi:10.1007/s10107-024-02164-2}.
\bibitem[DT14]{DT}
Y. Drori and M. Teboulle.
Performance of First-Order Methods for Smooth Convex Minimization: A Novel Approach.
\emph{Mathematical Programming} 145, 451--482, 2014.
\href{https://doi.org/10.1007/s10107-013-0653-0}{doi:10.1007/s10107-013-0653-0}.
\bibitem[HC26]{HC}
H. He and G. Cheng.
Optimal Polynomial Exponents of Gradient Descent with Arbitrary Real Stepsizes.
\href{https://ssrn.com/abstract=7569138}{SSRN 7569138}, 2026.
\bibitem[JCY26]{JCY}
M. Jung, H. Cho, and C. Yun.
Stronger Lower Bounds for (Non-)Anytime Acceleration of Gradient Descent.
\href{https://arxiv.org/abs/2609.04032v2}{arXiv:2609.04032v2}, 2026.
\bibitem[JH26]{JH}
Y. Jeong and Y. Hong.
Silver Rate Is Sharp for Gradient Descent.
\href{https://arxiv.org/abs/2611.30100v1}{arXiv:2611.30100v1}, 2026.
\bibitem[LY26]{LY}
K. Liu and Y. Ye.
Silver Rate Is (Almost) Optimal for Gradient Descent: The Strongly Convex Case.
\href{https://arxiv.org/abs/2609.26873v2}{arXiv:2609.26873v2}, 2026.
\bibitem[MC26]{MC}
J. Ma and Y. Chen.
A Lower Bound for Stepsize-Based Acceleration of Gradient Descent.
\href{https://arxiv.org/abs/2608.10418v1}{arXiv:2608.10418v1}, 2026.
\bibitem[Mor65]{Moreau}
J.-J. Moreau.
Proximit\'e et dualit\'e dans un espace hilbertien.
\emph{Bulletin de la Soci\'et\'e Math\'ematique de France} 93, 273--299, 1965.
\href{https://doi.org/10.24033/bsmf.1625}{doi:10.24033/bsmf.1625}.
\bibitem[THG17]{THG}
A.~B. Taylor, J.~M. Hendrickx and F. Glineur.
Smooth Strongly Convex Interpolation and Exact Worst-Case Performance of First-Order Methods.
\emph{Mathematical Programming} 161, 307--345, 2017.
\href{https://doi.org/10.1007/s10107-016-1009-3}{doi:10.1007/s10107-016-1009-3}.
\bibitem[YL26]{YLsmooth}
Y. Ye and K. Liu.
Silver Rate Is (Almost) Optimal for Gradient Descent.
\href{https://arxiv.org/abs/2609.09152v3}{arXiv:2609.09152v3}, 2026.
\end{thebibliography}
\end{document}
